\documentclass[
	% -- opções da classe memoir --
	article,			% indica que é um artigo acadêmico
	11pt,				% tamanho da fonte
	oneside,			% para impressão apenas no recto. Oposto a twoside
	a4paper,			% tamanho do papel. 
	twocolumn,
	% -- opções da classe abntex2 --
	%chapter=TITLE,		% títulos de capítulos convertidos em letras maiúsculas
	%section=TITLE,		% títulos de seções convertidos em letras maiúsculas
	%subsection=TITLE,	% títulos de subseções convertidos em letras maiúsculas
	%subsubsection=TITLE % títulos de subsubseções convertidos em letras maiúsculas
	% -- opções do pacote babel --
	english,			% idioma adicional para hifenização
	brazil,				% o último idioma é o principal do documento
	sumario=tradicional
	]{abntex2}

\usepackage{cmap}	
\usepackage{lmodern}	
\usepackage[T1]{fontenc}	
\usepackage[utf8]{inputenc}		
\usepackage{lastpage}		
\usepackage{color}	
\usepackage{graphicx}	
\usepackage{units}
\usepackage[brazilian,hyperpageref]{backref}
\usepackage[num]{abntex2cite}
\usepackage{bold-extra}
\usepackage{eso-pic}
\usepackage{indentfirst, microtype}		
\usepackage{caption,subcaption,float, dcolumn}
\usepackage{amsfonts,amsmath,amssymb}
\usepackage{enumerate}
\usepackage{titlesec}

\usepackage{multirow}


\usepackage{algpseudocode}
\usepackage{algorithm}

\floatname{algorithm}{Algoritmo}
\renewcommand{\algorithmicrequire}{\textbf{Input:}}
\renewcommand{\algorithmicensure}{\textbf{Output:}}

\usepackage{svg}

\usepackage{listings}
\usepackage[newfloat]{minted}
% \usemintedstyle{friendly}
\citebrackets[]

\usemintedstyle{friendly}
\newenvironment{code}{\captionsetup{type=listing}}{}
\SetupFloatingEnvironment{listing}{name=Código}

\newcommand{\algorithmautorefname}{Algoritmo}
\renewcommand{\backrefpagesname}{Citado na(s) página(s):~}
\renewcommand{\backref}{}
\renewcommand*{\backrefalt}[4]{
	\ifcase #1 %
		Nenhuma citação no texto.%
	\or
		Citado na página #2.%
	\else
		Citado #1 vezes nas páginas #2.%
	\fi}%
% ---

\definecolor{blue}{rgb}{0, .226, .437}
\makeatletter
\hypersetup{
pdftitle={Circuito digital CMOS para controle do fator de qualidade de um filtro passa-banda ativo sintonizável}, 
		pdfauthor={Alef de Oliveira Santos},
    	pdfsubject={TCC2},
	    pdfcreator={LaTeX with abnTeX2},
        pdfkeywords={Fator Q. Controle digital. Métodos numéricos. RTL. Verilog.},
		colorlinks=true,       		% false: boxed links; true: colored links
    	linkcolor=blue,          	% color of internal links
    	citecolor=blue,        		% color of links to bibliography
    	filecolor=magenta,      		% color of file links
		urlcolor=blue,
		bookmarksdepth=4
}
\makeatother
\setlength{\parindent}{1.3cm}
\setlength{\parskip}{0.2cm}  
\makeindex

\renewcommand{\listalgorithmname}{Lista de algoritmos}

\titleformat{\chapter}     {\normalfont\normalsize\bfseries}{\thechapter}     				 {1em}{}
\titleformat{\section}     {\normalfont\normalsize}{\thechapter.\thesection}     			 {1em}{}
\titleformat{\subsection}  {\normalfont\normalsize}{\thechapter.\thesection.\thesubsection}  {1em}{}


\setlrmarginsandblock{1.5cm}{1.5cm}{*}
\setulmarginsandblock{1.5cm}{2.5cm}{*}
\checkandfixthelayout




\makeindex
% ---

% ---
% Altera as margens padrões
% ---
\titleformat{\chapter}     {\normalfont\normalsize\bfseries}{\thechapter}     				 {1em}{}
\titleformat{\section}     {\normalfont\normalsize}{\thechapter.\thesection}     			 {1em}{}
\titleformat{\subsection}  {\normalfont\normalsize}{\thechapter.\thesection.\thesubsection}  {1em}{}
% ---

% --- 
% Espaçamentos entre linhas e parágrafos 
% --- 

% O tamanho do parágrafo é dado por:
\setlength{\parindent}{1.3cm}

% Controle do espaçamento entre um parágrafo e outro:
\setlength{\parskip}{0.2cm}  % tente também \onelineskip

% Espaçamento simples
\SingleSpacing


% ----
% Início do documento
% ----
\begin{document}

% Seleciona o idioma do documento (conforme pacotes do babel)
%\selectlanguage{english}
\selectlanguage{brazil}

% Retira espaço extra obsoleto entre as frases.
\frenchspacing 

% ----------------------------------------------------------
% ELEMENTOS PRÉ-TEXTUAIS
% ----------------------------------------------------------

%---
%
% Se desejar escrever o artigo em duas colunas, descomente a linha abaixo
% e a linha com o texto ``FIM DE ARTIGO EM DUAS COLUNAS''.
% \twocolumn[    		% INICIO DE ARTIGO EM DUAS COLUNAS
%
%---

% página de titulo principal (obrigatório)
% \maketitle
% titulo em outro idioma (opcional)

\title{Circuito digital CMOS para controle do fator de qualidade de um filtro passa-banda ativo sintonizável}

\twocolumn[
	\begin{center}
		% % \begin{minipage}[H]{0.4}
		% % 	\begin{figure*}[]
		% % 		\centering
		% % 		\includegraphics[H]{./unifeilogo.eps}
				
		% % 	\end{figure*}
			
		% % \end{minipage}
		% \begin{minipage}
			
		% \end{minipage}
		\Large 
		 	Trabalho de conclusão de curso \\
		 	Universidade Federal de Itajubá - \textit{Campus} Theodomiro Carneiro Santiago\\
			Instituto de Ciências Tecnológicas - ICT\\
			\large
			Engenharia de Computação
		
		\normalsize
			
			
			\vspace*{1cm}
			\textbf{\thetitle}
			\vspace*{1cm}

			Alef de Oliveira Santos${}^*$, Dean Bicudo Karolak${}^*$, Paulo Márcio Moreira e Silva${}^*$

			\vspace*{0.5cm}

			${}^*$\textit{Universidade Federal de Itajubá - Campus Theodomiro Carneiro Santiago \\ Rua Rua Irmã Ivone Drumond, 200 - Distrito Industrial II - 35903-087 \\
			Itabira, Minas Gerais, Brasil} \\

			\vspace*{0.5cm}

			\textit{E-mail}: 
				\href{mailto:alef_santos@unifei.edu.br}{\verb|alef\_santos@unifei.edu.br}, 
				\href{mailto:dean.karolak@unifei.edu.br}{\verb|dean.karolak@unifei.edu.br},
				\href{mailto:paulo.silva@unifei.edu.br }{\verb|paulo.silva@unifei.edu.br } 
		
		
	\end{center}

	\noindent

	\textbf{Abstract:}\\

	\textbf{Keywords:}\\

	\textbf{Resumo:} 
	O controle do fator de qualidade ($Q$) em um filtro passa banda é uma maneira direta de controlar a seletividade deste filtro, e utilizando um circuito ativo é possível controlá-lo. Para controlar o $Q$ de um sistema ressonante com um circuito digital utiliza-se neste trabalho técnicas de computação numérica para aproximação de funções não-lineares, uma vez que os parâmetros do filtro \textit{on-chip} não são manipuláveis além da corrente que controla o $Q$. Assim, o objetivo deste trabalho é projetar um circuito digital sintetizável capaz de: receber como entrada o $Q$  desejado e medido de um filtro passa banda e controlar digitalmente esse valor de $Q$ do filtro. Além disso, três diferentes métodos numéricos de controle serão estudados e implementados para comparação e escolha do melhor método de convergência. \\
	\textbf{Palavras-Chave:}\\
]


% resumo em português
% \begin{resumoumacoluna}
%  Conforme a ABNT NBR 6022:2018, o resumo no idioma do documento é elemento obrigatório. 
%  Constituído de uma sequência de frases concisas e objetivas e não de uma 
%  simples enumeração de tópicos, não ultrapassando 250 palavras, seguido, logo 
%  abaixo, das palavras representativas do conteúdo do trabalho, isto é, 
%  palavras-chave e/ou descritores, conforme a NBR 6028. (\ldots) As 
%  palavras-chave devem figurar logo abaixo do resumo, antecedidas da expressão 
%  Palavras-chave:, separadas entre si por ponto e finalizadas também por ponto.
 
%  \vspace{\onelineskip}
 
%  \noindent
%  \textbf{Palavras-chave}: latex. abntex. editoração de texto.
% \end{resumoumacoluna}


% % resumo em inglês
% \renewcommand{\resumoname}{Abstract}
% \begin{resumoumacoluna}
%  \begin{otherlanguage*}{english}
%    According to ABNT NBR 6022:2018, an abstract in foreign language is optional.

%    \vspace{\onelineskip}
 
%    \noindent
%    \textbf{Keywords}: latex. abntex.
%  \end{otherlanguage*}  
% \end{resumoumacoluna}

% % ]  				% FIM DE ARTIGO EM DUAS COLUNAS
% % ---

% \begin{center}\smaller
% \textbf{Data de submissão e aprovação}: elemento obrigatório. Indicar dia, mês e ano

% \textbf{Identificação e disponibilidade}: elemento opcional. Pode ser indicado 
% o endereço eletrônico, DOI, suportes e outras informações relativas ao acesso.
% \end{center}

% ----------------------------------------------------------
% ELEMENTOS TEXTUAIS
% ----------------------------------------------------------
\textual

% ----------------------------------------------------------
% Introdução
% ----------------------------------------------------------
\chapter{Introdução}

% o que é o quê e por quê ele é importante?
% como medir o Q
% [ENFASE] como controlar o Q 
% 

Este trabalho trata de um sistema eletrônico que recebe e controla o fator de qualidade ($Q$) de um circuito eletrônico ressonante. O circuito proposto utiliza-se de técnicas de computação numérica para controlar o fator de qualidade através de uma corrente de referência injetada no sistema ressonante. Neste trabalho serão realizados circuitos digitais periféricos para a determinação do valor de $Q$ medido de um filtro passa banda ativo, bem como, um circuito digital para controle e aproximação do $Q$ desejado. Em relação ao controle e aproximação de Q, serão comparados os métodos numéricos da Bisseção, Secantes e Secantes com seleção de intervalo implementados. O sistema digital é projetado e implementado em Verilog tendo em mente a posterior fabricação em silício. Para efeitos de estudo e desenvolvimento, este projeto em ASIC utiliza uma tecnologia GPDK de 45nm.
% ----------------------------------------------------------
% ELEMENTOS PÓS-TEXTUAIS
% ----------------------------------------------------------
\postextual

% ----------------------------------------------------------
% Referências bibliográficas
% ----------------------------------------------------------
\bibliography{ref}

\end{document}
